% apaquasi: the first page of a preprint, set the way a high-impact journal
% sets the first page of an article. The model is Nature Human Behaviour.
%
% Loaded after apalatex.tex, and used only by the front matter quasifront.lua
% builds in journal mode. Everything after the first page is apaquarto's.
%
% The page, from the top:
%
%   * a thick bar in the accent colour, and under it the article type as a
%     wordmark (titlepage: type:, Research Article by default), the way a
%     journal prints the section an article belongs to;
%   * a label at the left (titlepage: label:, Preprint by default) and the doi
%     at the right, over a hairline;
%   * the title: very large, heavy, tightly set serif;
%   * a hairline across the page, and under it two columns: a sidebar of
%     entries ruled off from one another (date, version, status, keywords,
%     materials, open science badges) at the left; the byline with the
%     orcids, the affiliations, a hairline and the abstract at the right;
%   * the author note and the correspondence address at the foot of the first
%     column, and the article's line at the foot of the page.

\usepackage{xcolor}
\usepackage{marvosym}
% The icons of the materials list: fontawesome5 (github, gitlab, a link) and
% academicons (osf, zenodo, psyarxiv, figshare). Both ship with TeX Live.
\usepackage{fontawesome5}
\usepackage{academicons}

% The one colour the page carries: the bar, the doi and the foot line. A
% document sets its own with titlepage: color: "B8202E" (hex, no hash).
\definecolor{quasiaccent}{HTML}{1B5E9F}
\definecolor{quasigrey}{gray}{0.35}

% The line numbers in the margin, there for reviewers to point at a line: a
% light grey, so that they sit back from the text. lineno is loaded by
% formatlatex.lua, in the header too, and the order of the two is not
% promised, so the font is set once the document begins, and only if lineno
% is there to use it. Its own size and face, lighter.
\definecolor{quasilinenumber}{gray}{0.62}
\AtBeginDocument{%
  \ifdefined\linenumberfont
    \renewcommand{\linenumberfont}{\normalfont\tiny\sffamily\color{quasilinenumber}}%
  \fi}

% --- Faces --------------------------------------------------------------------
% The title is set in a heavy serif, as the Nature journals set theirs: TeX
% Gyre Bonum Bold (Bookman), which ships with every TeX distribution alongside
% Termes and Heros. Roboto Serif was used here first, but its TeX Live build
% keeps the overlapping outlines of the variable font it was cut from (the
% stem and the bar of a t, say), which Acrobat draws with seams at some zooms.
\newfontfamily\quasititlefamily{texgyrebonum-bold.otf}[LetterSpace=-2]
\newfontfamily\quasiwordmarkfamily{texgyrebonum-bold.otf}[LetterSpace=-2]

% Hairlines: the rules the page is divided up with.
\newlength{\quasihairline}
\setlength{\quasihairline}{0.4pt}
\newcommand{\quasirule}{%
  \par\nointerlineskip\noindent\rule{\linewidth}{\quasihairline}\par\nointerlineskip}

% --- The height of it all ------------------------------------------------------
% Everything above the two columns is one box, which \twocolumn sets across the
% top of the first page, and a box does not break: an abstract far past APA's
% 250 words, a long sidebar or a long list of authors made it taller than the
% page, and it ran into the foot margin and pushed the article to page 2. A box
% taller than 0.88 of the text height is scaled down to that, centred, which
% leaves room for the first lines of the article, and latex says so in the log.
% apalatex.tex's \apamastheadlift is what \twocolumn calls just before it sets
% the box, so the measuring goes there.
\newlength{\quasimastheadmax}
\newcommand{\quasifitmasthead}{%
  \setlength{\quasimastheadmax}{0.88\textheight}%
  \ifdim\dimexpr\ht\apamastheadbox+\dp\apamastheadbox\relax>\quasimastheadmax
    \PackageWarningNoLine{apaquasi}{The front matter of the first page is taller
      than the page allows,\MessageBreak and has been scaled down to fit. Shorten
      the abstract,\MessageBreak the title or the sidebar to set it at full size}%
    \sbox{\apamastheadbox}{\makebox[\textwidth]{%
      \resizebox*{!}{\quasimastheadmax}{\usebox{\apamastheadbox}}}}%
  \fi}
\let\quasiapamastheadlift\apamastheadlift
\renewcommand{\apamastheadlift}{\quasifitmasthead\quasiapamastheadlift}

% --- The band -----------------------------------------------------------------
\newcommand{\quasibar}{%
  \noindent{\color{quasiaccent}\rule{\linewidth}{4pt}}\par\nointerlineskip}

% The wordmark under the bar, when there is one.
\newcommand{\quasiwordmark}[1]{%
  \vspace{4pt}\noindent{\quasiwordmarkfamily\fontsize{20}{24}\selectfont #1}\par}

% The label at the left and the doi at the right, over a hairline.
\newcommand{\quasitypeline}[2]{%
  \vspace{22pt}\noindent
  {\sffamily\bfseries\fontsize{11}{13}\selectfont #1}%
  \hfill{\sffamily\fontsize{9}{11}\selectfont\color{quasiaccent}#2}%
  \par\vspace{3pt}\quasirule\vspace{6pt}}

% --- Title --------------------------------------------------------------------
\newcommand{\quasititlefont}{\quasititlefamily\fontsize{26}{29}\selectfont}

\newenvironment{quasititle}
  {\par\noindent\begin{minipage}{\linewidth}\raggedright\quasititlefont}
  {\end{minipage}\par}

% --- The two columns under the title ------------------------------------------
% The sidebar takes under a third of the measure, the byline and the abstract
% the rest, less a gutter. Measured where they are used: in the preamble,
% \linewidth is not yet the page's.
\newlength{\quasisidewidth}
\newlength{\quasimainwidth}

\newcommand{\quasisidebar}{%
  \setlength{\quasisidewidth}{0.31\linewidth}%
  \setlength{\quasimainwidth}{0.65\linewidth}%
  \par\vspace{26pt}\quasirule\vspace{6pt}\noindent
  \begin{minipage}[t]{\quasisidewidth}%
  \vspace{0pt}%
  \sffamily\fontsize{8.5}{11}\selectfont
  \setlength{\parindent}{0pt}\setlength{\parskip}{0pt}\raggedright}

% One entry of the sidebar, ruled off from the next.
\newcommand{\quasisidelabel}[1]{#1: }
\newcommand{\quasisiderule}{\par\vspace{4pt}\quasirule\vspace{5pt}}

% An academicons icon. The package switches to its font without a group of
% its own, which left the text after an icon in a font with no letters in it.
\newcommand{\quasiai}[1]{{\aiicon{#1}}}

% One line of the materials list: the icon of the repository as its bullet,
% and the link, hanging clear of the bullet when it wraps.
\newlength{\quasibulletwidth}
\setlength{\quasibulletwidth}{1.35em}
\newcommand{\quasimaterial}[2]{%
  \par\noindent\hangindent=\quasibulletwidth\hangafter=1
  \makebox[\quasibulletwidth][l]{\fontsize{9.5}{11}\selectfont#1}#2}

% The open science badges: the Center for Open Science's own files, which
% carry their names on their ribbons (badges/, from https://osf.io/tvyxz/).
% Three to a row of the sidebar, or four, a little smaller, when there are
% more than three: quasifront.lua works out the width, as a fraction of the
% sidebar, and passes it here. The badges are a self-disclosure (see the
% README).
\newlength{\quasibadgewidth}
\newcommand{\quasibadge}[1]{%
  \includegraphics[width=\quasibadgewidth]{#1}}
\newcommand{\quasibadgegap}{\hspace{0.035\linewidth}\allowbreak}
% The row, with a small heading of its own.
\newcommand{\quasibadges}[1]{%
  \setlength{\quasibadgewidth}{#1\linewidth}%
  \noindent Open science\par\vspace{5pt}%
  \noindent\setlength{\lineskip}{5pt}\setlength{\lineskiplimit}{5pt}}

% The byline and the affiliations under it, over a hairline of their own, then
% the abstract.
\newcommand{\quasimain}{%
  \end{minipage}\hfill
  \begin{minipage}[t]{\quasimainwidth}%
  \vspace{0pt}%
  \raggedright\sffamily\bfseries\fontsize{10.5}{13}\selectfont
  \setlength{\parindent}{0pt}}

\newcommand{\quasiaffiliations}{%
  \par\vspace{5pt}%
  \mdseries\fontsize{8}{10}\selectfont\color{quasigrey}}

\newcommand{\quasiabstract}{%
  \par\vspace{18pt}\quasirule\vspace{6pt}%
  \color{black}\rmfamily\mdseries\fontsize{11}{15}\selectfont
  \setlength{\parindent}{0pt}\setlength{\parskip}{5pt}%
  \setlength{\RaggedRightParindent}{0pt}\RaggedRight
  % Ragged right, and so rarely hyphenated.
  \hyphenpenalty=5000\exhyphenpenalty=5000}

\newcommand{\quasimainend}{%
  \end{minipage}%
  \par\vspace{16pt}}

% The envelope that marks the corresponding author, and heads the address at
% the foot of the page.
\newcommand{\quasienvelope}{\raisebox{0.05ex}{\scriptsize\Letter}}

% The journals set the first paragraph of the article flush left. Asked for
% once, before the body: the next paragraph to start throws its indent away.
\newcommand{\quasifirstpar}{\everypar{{\setbox0\lastbox}\everypar{}}}

% --- The foot of the first column ---------------------------------------------
% The author note and the address, small and in the grotesque, under a
% hairline. apajounote, which apalatex.tex defines and formatlatex.lua raises,
% puts them at the foot of the column.
\renewcommand{\apajounotetext}{%
  \sffamily\fontsize{7}{9}\selectfont
  \setlength{\parindent}{0pt}\setlength{\parskip}{2pt}%
  \setlength{\RaggedRightParindent}{0pt}\RaggedRight}

\renewcommand{\apajounoterule}{%
  \noindent\rule{\linewidth}{\quasihairline}%
  \par\vspace{3pt}}

% The author note is set there too, which makes the float taller than latex
% will hold at the foot of this column. Latex allows a bottom float half the
% column, and on the first page the column is only what the masthead leaves of
% it, so the note went to the second page. For the note alone, and only until
% the column ends, the foot of the column may take nine tenths of it.
\makeatletter
\AddToHook{env/apajounote/before}{\global\@botroom=0.9\@colht}
\makeatother

% --- The foot of the first page -----------------------------------------------
% The journals print their name at the foot of the opening page, in the accent
% colour. A preprint prints its own line there, out of what it has. Global,
% since the masthead that sets it is built in a box.
\newcommand{\quasifootline}{}
\newcommand{\setquasifootline}[1]{\gdef\quasifootline{#1}}

\fancypagestyle{apajoufirstpage}{%
  \fancyhf{}%
  \renewcommand{\headrulewidth}{0pt}%
  \fancyfoot[L]{{\sffamily\fontsize{8}{10}\selectfont\color{quasiaccent}%
    \quasifootline}}%
  \fancyfoot[R]{{\sffamily\fontsize{8}{10}\selectfont\thepage}}}
